% AI Explain Working Paper: LaTeX template.
%
% Google Scholar reads the first page of the PDF, so keep this layout: the title in
% a large font at the top, the authors directly below it, the abstract, and a
% section headed "References" at the end. Export a text-based PDF (pdfLaTeX,
% XeLaTeX, LuaLaTeX and Overleaf all do) and keep it under 5 MB.

\documentclass[11pt,a4paper]{article}

\usepackage[T1]{fontenc}
\usepackage{lmodern}
\usepackage[margin=25mm]{geometry}
\usepackage{amsmath,amssymb}
\usepackage{graphicx}
\usepackage{fancyhdr}
\usepackage[hidelinks]{hyperref}

% Paper details: keep them in sync with papers/<id>/paper.yaml.
\newcommand{\paperid}{AIE-WP-2026-001}
\newcommand{\paperversion}{v1.0}
\newcommand{\paperdate}{26 October 2026}
\newcommand{\papertitle}{Title of the Paper}

\hypersetup{
  pdftitle={\papertitle},
  pdfauthor={First Author, Second Author},
  pdfsubject={AI Explain Working Paper \paperid},
}

\setlength{\headheight}{14pt}
\pagestyle{fancy}
\fancyhf{}
\renewcommand{\headrulewidth}{0.4pt}
\fancyhead[L]{\small AI Explain Working Paper \paperid}
\fancyhead[R]{\small Version \paperversion, \paperdate}
\fancyfoot[C]{\thepage}

\title{\papertitle}
\author{
  First Author\\ \small AI Explain
  \and
  Second Author\\ \small Affiliation
}
\date{}

\begin{document}

\maketitle
\thispagestyle{fancy}

\begin{abstract}
Write the abstract here. Use the same text as the abstract in paper.yaml.
\end{abstract}

\noindent\textbf{Keywords:} first keyword, second keyword

\section{Introduction}

Start the paper here and cite related work like this~\cite{fips180}.

\section{Conclusion}

Summarise the findings.

% The article class titles this section "References".
\begin{thebibliography}{9}

\bibitem{fips180}
National Institute of Standards and Technology.
\newblock \emph{Secure Hash Standard (SHS)}.
\newblock FIPS PUB 180-4, 2015.

\end{thebibliography}

\end{document}
